\chapter{固有値,固有ベクトル}
    \section{諸定理}
        \subsection{逆行列の特性多項式: \texorpdfstring{$\phi_{A^{-1}}(\lambda) = |A|^{-1}(-\lambda)^n\phi_A(1/\lambda)$}{φA\string^(-1)(λ) = |A|\string^(-1)\string^n φA(1/λ)}}
            \begin{shadebox}
                $A$を正則な$n$次正方行列とし，特性多項式を$\phi_A(\lambda)$とすると，$\phi_{A^{-1}}(\lambda) = |A|^{-1}(-\lambda)^n\phi_A(1/\lambda)$。
            \end{shadebox}
            \begin{proof}
                \quad\par
                $|A(\lambda I - A^{-1})|$を2通りの方法で評価する。
                まず$|A(\lambda I - A^{-1})| = |A|\phi_{A^{-1}}(\lambda)$である。
                次に$|A(\lambda I - A^{-1})| = |\lambda A - I| = |(-\lambda)((1/\lambda)I - A)| = (-\lambda)^n|(1/\lambda)I - A| = (-\lambda)^n\phi_A(1/\lambda)$。
                よって主張が成り立つ。
            \end{proof}
        \subsection{相異なる固有値に対応する固有ベクトルは一次独立}
            \label{相異なる固有値に対応する固有ベクトルは一次独立}
            \begin{shadebox}
                $A\in\complexNumbers^{n\times n}$の相異なる$p$個の固有値$\lambda_i\;(i=1,\dotsc,p)$に対応する固有ベクトルを$\bm{v}_i$とすると$\bm{v}_1,\dotsc,\bm{v}_p$は一次独立である。
            \end{shadebox}
            \begin{proof}
                \quad\par
                $c_1,\dotsc,c_p$を定数として
                \[c_1\bm{v}_1 + \dotsb + c_p\bm{v}_p = \bm{0}_n\]
                とおいて両辺に左から$A$を掛けると
                \[\lambda_1 c_1\bm{v}_1 + \dotsb + \lambda_p c_p\bm{v}_p = \bm{0}_n\]
                となる。これを繰り返して
                \begin{align*}
                    {\lambda_1}^2 c_1\bm{v}_1 + \dotsb + {\lambda_p}^2 c_p\bm{v}_p &= \bm{0}_n \\
                    & \vdots \\
                    {\lambda_1}^{p-1} c_1\bm{v}_1 + \dotsb + {\lambda_p}^{p-1} c_p\bm{v}_p &= \bm{0}_n
                \end{align*}
                を得る。これより
                \[
                    [c_1\bm{v}_1,\dotsc,c_p\bm{v}_p]
                    \begin{bmatrix}
                        1 & \lambda_1 & \cdots & {\lambda_1}^{p-1} \\
                        1 & \lambda_2 & \cdots & {\lambda_2}^{p-1} \\
                        & \vdots & \\
                        1 & \lambda_p & \cdots & {\lambda_p}^{p-1}
                    \end{bmatrix}
                    = O_{n\times p}
                \]
                となるが，$\lambda_i\;(i=1,\dotsc,p)$が相異なることより，上式の等号の左隣の Vandermonde 行列は正則だから
                \[ [c_1\bm{v}_1,\dotsc,c_p\bm{v}_p] = O_{n\times p} \]
                となる。$\bm{v}_1,\dotsc,\bm{v}_p\neq\bm{0}_n$より$c_1=\dotsb=c_p=0$である。
            \end{proof}
        \subsection{\texorpdfstring{$n\times n$}{n×n}行列\texorpdfstring{$A$}{A}の全要素が\texorpdfstring{$a\neq 0$}{a≠0}であるなら非零固有値は\texorpdfstring{$an$}{an}のみ}
            \begin{proof}
                \quad\par
                まず，$an$が固有値であり，全ての要素が等しい$n$次元ベクトルが固有値$an$に対応する固有ベクトルであることは，実際に計算してみることで容易に確かめられる。
                非零固有値がこれ以外に存在しないことを示す。$\lambda\neq 0,an$とし，$\lambda I-A$の第$j$列ベクトルを$\bm{b}_j$とすると，
                \[
                    b_j(i) =
                    \begin{cases}
                        -a & (i\neq k)\\
                        \lambda-a & (i=k)
                    \end{cases}
                \]
                $\{\bm{b}_1,\dotsc,\bm{b}_n\}$が一次独立であることを示せば，$\left|\lambda I-A\right|\neq 0$が示せて，$\lambda$が固有値でないことが示せる。
                $c_1,\dotsc,c_n$を定数とし，$\sum_{j=1}^{n}{c_j\bm{b}_j}=\bm{0}$とおくと，第$i$成分に着目して
                \[0 = c_i(\lambda-a) - a\sum_{j\neq i}^{n}{c_j} = c_i\lambda - a\sum_{j=1}^{n}{c_j}\]
                これと$\lambda\neq 0$より$c_1=\dotsb=c_n$である。この値を$c$とおくと$c\lambda = anc$となり，$\lambda\neq an$より$c=0$である。
            \end{proof}
        \subsection{転置行列の固有値,固有ベクトル}
            \begin{shadebox}
                $A\in\field^{n\times n}$について$A^\top$の固有値は$A$と等しく，各固有値$\lambda_i$に対する固有空間の次元も$A$の固有値$\lambda_i$に対するそれと等しい。
            \end{shadebox}
            \begin{proof}
                \quad\par
                $A^\top$の特性多項式は
                \[|sI-A^\top| = |(sI-A)^\top| = |sI-A|\]
                の如く$A$の特性多項式と等しいから固有値も全て等しい。
                各固有値$\lambda_i$に対する固有空間$E_{A^\top}(\lambda_i)$の次元は
                \begin{align*}
                    \dim E_{A^\top}(\lambda_i) &= \dim\Ker{\lambda_i I-A^\top} = \dim \field^n-\rank{\lambda_i I-A^\top} = n-\rank{(\lambda_i I-A)^\top}\\
                    &=n-\rank{\lambda_i I-A} = \dim E_A(\lambda_i)
                \end{align*}
                の如く$A$の固有値$\lambda_i$に対する固有空間の次元と等しい。
            \end{proof}
        \subsection{逆行列の固有値,固有ベクトル}
            \begin{shadebox}
                正則な行列$A\in\field^{n\times n}$の逆行列の固有値は$A$の固有値の逆数であり，$A$と共通の固有ベクトルをもつ。
            \end{shadebox}
            \begin{proof}
                \quad\par
                $A^{-1}$もまた正則だから0固有値は持たない。
                $A^{-1}$の固有値の一つを$\lambda'\neq0$とし，それに対する固有ベクトルを$\bm{y}\neq\bm{0}$とおくと
                $A^{-1}\bm{y} = \lambda'\bm{y}$であり，両辺に左から$A$を掛けて$\bm{y} = \lambda'A\bm{y}\quad\therefore\;A\bm{y} = \frac{1}{\lambda'}\bm{y}$となる。
                よって$A^{-1}$の固有値$\lambda'$は$A$の固有値$1/\lambda'$の逆数であり，固有ベクトルは共通である。
            \end{proof}
        \subsection{ユニタリ行列の全ての固有値の絶対値は1である}
            \begin{proof}
                \quad\par
                ユニタリ行列$P$とその任意の固有値$\lambda$およびそれに対応する固有ベクトル$\bm{x} \neq \bm{0}$について$P\bm{x} = \lambda\bm{x}$の両辺のノルムをとると$\|\bm{x}\|_2 = \|P\bm{x}\|_2 = |\lambda|\|\bm{x}\|_2$であるからわかる。
            \end{proof}
        \subsection{随伴行列}
            \subsubsection{定義}
                \label{随伴行列の定義}
                $n\in\naturalNumbers$ とし， $x$ を不定元とする。
                多項式 $f_n(x)\coloneq x^n+a_{n-1}x^{n-1}+\dotsb+a_0$ に対して次の行列
                \[
                    A =
                    \begin{bmatrix}
                        -a_{n-1} & -a_{n-2} & -a_{n-3} & \cdots & -a_0 \\
                        1 & 0 & 0 & \cdots & 0 \\
                        0 & 1 & 0 & \cdots & 0 \\
                        \vdots & & \ddots & \ddots \\
                        0 & 0 & 0 & 1 & 0
                    \end{bmatrix}
                \]
                を $f$ の\textbf{随伴行列}という。
            \subsubsection{行列式と固有多項式}
                \label{随伴行列の行列式と固有多項式}
                \begin{shadebox}
                    \ref{随伴行列の定義} の定義を引き継ぐ。
                    $\lambda\in\complexNumbers$ とする。
                    $\det{A}=(-1)^n a_0$ であり，$\det{(\lambda I-A)} = f_n(\lambda)$ である。
                    よって $A$ の固有値は $f_n$ の根である。
                \end{shadebox}
                \begin{proof}
                    \quad\par
                    $\det{A}=(-1)^n a_0$については，第1行をバブルソートの要領で$n-1$回の置換により最下行に移動させて下三角行列を作ることで証明できる。
                    \par
                    次に$\det{(\lambda I-A)} = f(\lambda)$を示す。
                    $n=1$のとき成り立つ。$n-1$まで成り立つと仮定する。
                    第1列に関して余因子展開すると
                    \begin{align*}
                        \det{(\lambda I-A)} &= \det
                        \begin{bmatrix}
                            \lambda+a_{n-1} & a_{n-2} & a_{n-3} & \cdots & a_0 \\
                            -1 & \lambda & 0 & \cdots & 0 \\
                            0 & -1 & \lambda & \cdots & 0 \\
                            \vdots & & \ddots & \ddots \\
                            0 & \cdots & 0 & -1 & \lambda
                        \end{bmatrix} \\
                        &= (\lambda+a_{n-1})\lambda^{n-1} + \det
                        \begin{bmatrix}
                            a_{n-2} & a_{n-3} & a_{n-4} & \cdots & a_0 \\
                            -1 & \lambda & 0 & \cdots & 0 \\
                            0 & -1 & \lambda & \cdots & 0 \\
                            \vdots & & \ddots & \ddots \\
                            0 & \cdots & 0 & -1 & \lambda
                        \end{bmatrix} \\
                        &= \lambda^{n} + a_{n-1}\lambda^{n-1} + \left.f_{n-1}(\lambda)\right|_{a_{n-2}\to a_{n-2}-\lambda} \\
                        &= \lambda^{n} + a_{n-1}\lambda^{n-1} + \lambda^{n-1} + (a_{n-2}-\lambda)\lambda^{n-2} + a_{n-3}\lambda^{n-3} + \dotsb + a_1\lambda + a_0 \\
                        &= \lambda^{n} + a_{n-1}\lambda^{n-1} + a_{n-2}\lambda^{n-2} + \dotsb + a_1\lambda + a_0 = f_n(\lambda)
                    \end{align*}
                \end{proof}
            \subsubsection{固有ベクトル}
                \label{随伴行列の固有ベクトル}
                \begin{shadebox}
                    \ref{随伴行列の定義} の定義を引き継ぐ。
                    $\lambda$ は $A$ の固有値であるとする\footnote{$\lambda=0$ も含む。何故ならば $\bm{v}\neq\bm{0}$ が必ず成り立つから。}。
                    $\bm{v}\coloneq[\lambda^{n-1},\lambda^{n-2},\dots,\lambda,1]^\top$ は $A$ の固有ベクトルである。
                \end{shadebox}
                \begin{proof}
                    \quad\par
                    \ref{随伴行列の行列式と固有多項式} より $-a_{n-1}\lambda^{n-1}- a_{n-2}\lambda^{n-2} - \dots - a_1\lambda - a_0 = \lambda^n$ であるから，次式が成り立つ。
                    \[ A\bm{v} = [\lambda^n,\lambda^{n-1},\dots,\lambda^2,\lambda] = \lambda\bm{v} \]
                \end{proof}
            \subsubsection{対角化}
                \begin{shadebox}
                    \ref{随伴行列の定義} の定義を引き継ぐ。
                    $A$ の固有値が全て異なるとし，それらを $\lambda_i\;(i=0,1,\dots,n)$ とする。
                    $\lambda_i$ に対応する固有ベクトルを $\bm{v}_i$ とする。
                    $A$ は $\lambda_i\;(i=0,1,\dots,n)$ に関する Vandermonde 行列 $V$ を用いて次のように対角化できる。
                    \[
                        A = V\Lambda V^{-1}\quad\text{where}\quad\Lambda\coloneq\diag{\lambda_0,\lambda_1,\dots,\lambda_n},\quad V\coloneq\begin{bmatrix}
                            \lambda_0^{n-1} & \lambda_1^{n-1} & \cdots & \lambda_n^{n-1} \\
                            \lambda_0^{n-2} & \lambda_1^{n-2} & \cdots & \lambda_n^{n-2} \\
                            \vdots \\
                            \lambda_0 & \lambda_1 & \cdots & \lambda_n \\
                            1 & 1 & \cdots & 1
                        \end{bmatrix}
                    \]
                \end{shadebox}
                \begin{proof}
                    \quad\par
                    $\lambda_i\;(i=0,1,\dots,n)$ が全て異なるから $V$ は正則である。
                    \ref{随伴行列の固有ベクトル} より $AV = V\Lambda$ が成り立つ。
                    この式の両辺に $V^{-1}$ を右から掛けて $A = V\Lambda V^{-1}$ を得る。
                \end{proof}
        \subsection{巡回行列の固有値}
            \begin{shadebox}
                $N\in\naturalNumbers,\;a_0,a_1,\dots,a_{N-1}\in\complexNumbers,\;\bm{a}\coloneq[a_0,\dots,a_{N-1}]^\top$とする。
                $\bm{a}$の要素を$m\in\{0,1,\dots,N-1\}$だけ巡回シフトしたベクトル$\bm{a}_m$を次式で定義する。
                \[
                    \bm{a}_m \coloneq
                    \begin{cases}
                        \bm{a} & (m=0) \\
                        [a_{N-m},\dots,a_{N-1},a_0,a_1,\dots,a_{N-1-m}]^\top & (m=1,2,\dots,N-1)
                    \end{cases}
                \]
                $A\in\complexNumbers^{N\times N}$を次式で定義する。
                \[ A \coloneq [\bm{a}_0,\bm{a}_1,\dots,\bm{a}_{N-1}] \]
                このとき$A$は次のように対角化可能である。
                \[ A = \sqrt{N}W\diag{\DFT{\bm{a}}}\HerConj{W} \]
                ここに$W$は$N\times N$ DFT基底行列であり，第$(n,k)$要素は次式で定義される。
                \[ W_{n,k} \coloneq \frac{1}{\sqrt{N}}\exp i\frac{nk}{N}2\pi \]
                $\DFT{\bm{a}}$は$\bm{a}$のDFT(離散 Fourier 変換)であり，次式で定義される。
                \[ \DFT{\bm{a}} \coloneq \HerConj{W}\bm{a} \]
                すなわち第$k$要素は次式である。
                \[ \sum_{n=0}^{N-1}\overline{W_{k,n}}a_n \]
            \end{shadebox}
            \begin{proof}
                \quad\par
                \cite{信号処理備忘録}で述べられているDFTの性質を用いる。
                $\bm{e}_m\in\realNumbers^N$を第$m$要素が1でそれ以外は0のベクトルとする。
                $\bm{a}_m$は巡回畳み込みを用いて次式で表せる。
                \[ \bm{a}_m = \bm{a}\circledast\bm{e}_m \]
                巡回畳み込みのDFTの性質より次式が成り立つ。
                \[ \DFT{\bm{a}_m} = \sqrt{N}\DFT{\bm{a}}\HadamardProd\DFT{\bm{e}_m} = \sqrt{N}\DFT{\bm{a}}\HadamardProd \HerConj{W}[:,m] \]
                これより次式が成り立つ。
                \begin{align*}
                    \HerConj{W}A &= [\DFT{\bm{a}_0},\DFT{\bm{a}_1},\dots,\DFT{\bm{a}_{N-1}}] \\
                    &= \sqrt{N}[\DFT{\bm{a}}\HadamardProd \HerConj{W}[:,0],\DFT{\bm{a}}\HadamardProd \HerConj{W}[:,1],\dots,\DFT{\bm{a}}\HadamardProd \HerConj{W}[:,N-1]] \\
                    &= \sqrt{N}\diag{\DFT{\bm{a}}}\HerConj{W}
                \end{align*}
                両辺に左から$W$を掛け，$W$のユニタリ性$W\HerConj{W} = I$を用いて次式を得る。
                \[ A = \sqrt{N}W\diag{\DFT{\bm{a}}}\HerConj{W} \]
            \end{proof}
        \subsection{Hadamard行列の固有値}
            \begin{shadebox}
                $2^n\;(n\in\naturalNumbers)$次のHadamard行列の固有値は$2^{n/2}$と$-2^{n/2}$が$2^{n-1}$個ずつである\cite{A Note on the Eigenvectors of Hadamard Matrices of Order 2^n}。
            \end{shadebox}
            \begin{proof}
                \quad\par
                $2^n$次のHadamard行列を$H_n$と書く。$H_0 \coloneq [1]$と定義する。
                まず，$H_1$の固有値と固有ベクトルの組が$(\pm\sqrt{2},[1\pm\sqrt{2},1]^\top)$(複合同順)であることは容易に導ける。
                よって$n=1$のとき定理の主張が成り立つ。
                \par
                $n=m$のとき成り立つと仮定して$n=m+1$のとき成り立つことを示す。
                $H_m$の固有値$2^{m/2}$に対応する相異なる固有ベクトルを$\bm{u}_m^1,\dotsc,\bm{u}_m^{2^{m-1}}$とし，固有値$-2^{m/2}$に対応する相異なる固有ベクトルを$\bm{v}_m^1,\dotsc,\bm{v}_m^{2^{m-1}}$とする。
                このとき次のことが成り立つ。\\
                $i=1,\dotsc,2^{m-1}$に対して
                \begin{equation}
                    \begin{bmatrix}
                        \bm{u}_m^i \\
                        (-1+\sqrt{2})\bm{u}_m^i
                    \end{bmatrix},\;
                    \begin{bmatrix}
                        (1-\sqrt{2})\bm{v}_m^i \\
                        \bm{v}_m^i
                    \end{bmatrix}
                    \text{は}H_{m+1}\text{の固有値}2^{(m+1)/2}\text{に対応する固有ベクトルである}
                \end{equation}
                \begin{equation}
                    \begin{bmatrix}
                        (1-\sqrt{2})\bm{u}_m^i \\
                        \bm{u}_m^i
                    \end{bmatrix},\;
                    \begin{bmatrix}
                        \bm{v}_m^i \\
                        (-1+\sqrt{2})\bm{v}_m^i
                    \end{bmatrix}
                    \text{は}H_{m+1}\text{の固有値}-2^{(m+1)/2}\text{に対応する固有ベクトルである}
                \end{equation}
                (1)の主張の半分は次のようにして確かめられる。
                \begin{align*}
                    H_{m+1}
                    \begin{bmatrix}
                        \bm{u}_m^i \\
                        (-1+\sqrt{2})\bm{u}_m^i
                    \end{bmatrix} =
                    \begin{bmatrix}
                        H_m & H_m \\
                        H_m & -H_m
                    \end{bmatrix}
                    \begin{bmatrix}
                        \bm{u}_m^i \\
                        (-1+\sqrt{2})\bm{u}_m^i
                    \end{bmatrix} =
                    \begin{bmatrix}
                        \sqrt{2}H_m\bm{u}_m^i \\
                        (2-\sqrt{2})H_m\bm{u}_m^i
                    \end{bmatrix} = 2^\frac{m+1}{2}
                    \begin{bmatrix}
                        \bm{u}_m^i \\
                        (-1+\sqrt{2})\bm{u}_m^i
                    \end{bmatrix}
                \end{align*}
                他も同様に確かめられる。
                これら$2^{m+1}$個の固有ベクトルが一次独立であることを示せば，$n$が1増える毎に$H_n$の固有ベクトルの個数が倍になるから，定理の証明が済む。
                相異なる固有値に対応する固有ベクトルは一次独立(\ref{相異なる固有値に対応する固有ベクトルは一次独立})だから，(1)で言及した固有ベクトルの一次結合と(2)で言及した固有ベクトルの一次結合は一次独立である。
                あとは，(1)の固有ベクトル達，および(2)の固有ベクトル達がそれぞれ一次独立であることを示せばよい。
                (1)の固有ベクトル達が一次独立であることを示す。(2)についても同様に示せる。
                (1)の列ベクトル達の上半分(或いは下半分)は$H_m$の固有ベクトルであるから，相異なる固有値に対応する固有ベクトルが一次独立であることを再び用いて
                $
                    \begin{bmatrix}
                        \bm{u}_m^i \\
                        (-1+\sqrt{2})\bm{u}_m^i
                    \end{bmatrix}
                    \;(i=1,\dotsc,2^{m-1})
                $
                の一次結合と
                $
                    \begin{bmatrix}
                        (1-\sqrt{2})\bm{v}_m^i \\
                        \bm{v}_m^i
                    \end{bmatrix}
                    \;(i=1,\dotsc,2^{m-1})
                $
                の一次結合は一次独立である。
                また，$\bm{u}_m^1,\dotsc,\bm{u}_m^{2^{m-1}}$は$H_m$の固有ベクトルだから
                $
                    \begin{bmatrix}
                        \bm{u}_m^i \\
                        (-1+\sqrt{2})\bm{u}_m^i
                    \end{bmatrix}
                    \;(i=1,\dotsc,2^{m-1})
                $
                は一次独立である。
                \par
                以上より(1),(2)で言及した固有ベクトル達は一次独立である。
            \end{proof}
        \subsection{最小消去多項式が最小多項式と一致するベクトルの存在}
            \begin{shadebox}
                任意の正方行列$A\in\field^{n\times n}$につきその最小消去多項式$\psi_{A,\vx}(\lambda)$が$A$の最小多項式$\psi(\lambda)$と一致するベクトル$\vx\in\field^n$が存在する。
            \end{shadebox}
            \begin{proof}
                \quad\par
                \[\psi(\lambda) = \prod_{i=1}^{p}{(\lambda-\lambda_i)^{\mu_i}}\]
                とする。但し$\lambda_i$は$A$の相異なる$p$個の固有値である。このとき$\field^n$は次のように直和分解できる(\cite{新微分方程式対話}\;P77\;「$\NapierE^{Mt}$と一般固有値問題」)。
                \[\field^n = \bigoplus_{i=1}^p{\Ker{(A-\lambda I)^{\mu_i}}}\]
                $\psi(\lambda)$より低次のいかなる多項式によっても消去できない$\vx\in\field^n$を構成しよう。
                \[\Ker{(A-\lambda_i I)^{\mu_i}}/\Ker{(A-\lambda_i I)^{\mu_i-1}}\supset\{\bm{0}\}\]
                である(さもなくば$\psi(\lambda)$が最小多項式であったことに矛盾する!)から，ここからひとつベクトル$\vx_i$をとると自然数$k_i$につき
                \[
                    (A-\lambda I)^{k_i}\vx_i
                    \begin{cases}
                        =\bm{0} & (\lambda=\lambda_i,\;k_i\geq\mu_i)\\
                        \neq\bm{0} & (\mathrm{otherwise})
                    \end{cases}
                \]
                もっと簡単に言うと，$\vx_i$を消去するには$(A-\lambda_i I)$の$\mu_i$次以上の積が必須である。
                そこで
                \[x = \sum_{i=1}^{p}{\vx_i}\]
                とすれば$\vx$は$\psi(\lambda)$で消去でき，かつ$\vx$の任意の消去多項式は必ず$\psi(\lambda)$で割り切れなくてはならないから$\vx$の最小消去多項式は$\psi(\lambda)$である。
            \end{proof}
    \section{スペクトル写像定理}
        \begin{shadebox}
            $A\in\field^{n\times n}$の固有値を重複も含めて$\lambda_1,\dotsc,\lambda_n$とし，$f$を多項式とし， $s$ を不定元とする。
            このとき行列多項式$f(A)$の固有値は重複を含めて$f(\lambda_1),\dotsc,f(\lambda_n)$となる。
        \end{shadebox}
        \begin{proof}
            \quad\par
            $A$は正方行列であるから適当なユニタリ行列$P\in\complexNumbers^{n\times n}$により次のように上三角化できる。
            \[A = PUP^\top\]
            但し$U$は第$i$対角成分が$\lambda_i$である上三角行列である($P$はユニタリだから$A^n = PU^nP^\top$となることに注意せよ)。
            $f(s)$の次数を$m$とし，$k$次の係数を$c_k$とする($f(s) = \sum_{k=0}^{m}c_{k}s^k$)。
            $f(A)$の固有多項式は
            \begin{align*}
                \left|\lambda I - f(A)\right| &= \left|\lambda I - \sum_{k=0}^{m}c_{k}A^k\right| = \left|\lambda I - \sum_{k=0}^{m}c_{k}PU^kP^\top\right| = \left|P\left(\lambda I - \sum_{k=0}^{m}c_{k}U^k\right)P^\top\right| \\
                &= \left|\lambda I - \sum_{k=0}^{m}c_{k}U^k\right|
            \end{align*}
            定理\ref{上(下)三角行列のべき乗}より$U^k$は上三角行列であり，その第$i$対角成分は${\lambda_i}^k$であるから，上の$|\cdot|$の中身は上三角行列であり，その第$i$対角成分は$\lambda - \sum_{k=0}^{m}c_{k}{\lambda_i}^k = \lambda - f(\lambda_i)$である。
            よって上の式は
            \[\left|\lambda I - \sum_{k=0}^{m}c_{k}U^k\right| = \prod_{i=1}^{n}{(\lambda - f(\lambda_i))}\]
            であるから$f(A)$の固有値が$f(\lambda_1),\dotsc,f(\lambda_n)$であることがわかる。
        \end{proof}
    \section{対角化}
        \subsection{正規行列\texorpdfstring{$\iff$}{⇔}ユニタリ行列で対角化可能}
            \begin{shadebox}
                $n$を自然数とする。
                $A \in \complexNumbers^{n\times n}$が正規行列であるための必要十分条件は，ある対角行列$\Lambda \in \complexNumbers^{n \times n}$とユニタリ行列$Q \in \complexNumbers^{n \times n}$が存在して$A = Q\Lambda \HerConj{Q}$となることである。
            \end{shadebox}
            \begin{proof}
                \quad\par
                十分性は容易に示せるので省略し，必要性のみ示す。
                $A$を正規行列とする。
                $A$のSchur分解を$A = SU\HerConj{S}$とする。
                ここに$S,U \in \complexNumbers^{n \times n}$はそれぞれ適当なユニタリ行列, 上三角行列である。
                $A$が正規だから
                \begin{align*}
                    O = \HerConj{A}A - A\HerConj{A} &= S\HerConj{U}U\HerConj{S} - SU\HerConj{U}\HerConj{S} = S(\HerConj{U}U - U\HerConj{U})\HerConj{S} \\
                    \therefore\;\HerConj{U}U - U\HerConj{U} &= \HerConj{S}OS = O \\
                    \therefore\;\HerConj{U}U &= U\HerConj{U}
                \end{align*}
                これより，まず$\HerConj{U}U$と$U\HerConj{U}$の第1,1成分が等しいから
                \[ (\HerConj{U}U)_{1,1} = (U\HerConj{U})_{1,1} \]
                \[ \therefore |u_{1,1}|^2 = |u_{1,1}|^2 + |u_{1,2}|^2 + \dotsb + |u_{1,n}|^2 \]
                \[ u_{1,2} = u_{1,3} = \dotsb = u_{1,n} = 0 \]
                次に，$\HerConj{U}U$と$U\HerConj{U}$の第2,2成分が等しいから
                \[ |u_{1,2}|^2 + |u_{2,2}|^2 = |u_{2,2}|^2 + |u_{2,3}|^2 + \dotsb + |u_{2,n}|^2 \]
                これと先程の結果$|u_{1,2}| = 0$より
                \[ u_{2,3} = u_{2,4} = \dotsb = u_{2,n} = 0 \]
                同じ要領で続けていくと$U$の非対角要素が全て0であることがわかる。
            \end{proof}
    \section{一般固有空間}
        \subsection{一般固有空間の階数の頭打ち}
            \begin{shadebox}
                行列$A$とその固有値$\lambda$の$k$階一般固有空間$F_{A,k}(\lambda)$に対して，もしある自然数$l$が存在して$F_{A,l}(\lambda)=F_{A,l+1}(\lambda)$が成り立てば，任意の$m\geq l+2$に対して$F_{A,m}(\lambda)=F_{A,l}(\lambda)$。
            \end{shadebox}
            \begin{proof}
                \quad\par
                $m=l+1$に対して示す。それ以上の$m$に対しても帰納的に示せる。
                結論を否定すると，ある$\bm{v}\in F_{A,m}(\lambda)$が存在して$(A-\lambda I)^{l+2}\bm{v}=\bm{0}\land(A-\lambda I)^{l+1}\bm{v}\neq\bm{0}$。
                すなわち$(A-\lambda I)^{l+1}\left[(A-\lambda I)\bm{v}\right]=\bm{0}\land(A-\lambda I)^l\left[(A-\lambda I)\bm{v}\right]\neq\bm{0}$。
                よって$(A-\lambda I)\bm{v} \in F_{A,l+1}(\lambda)$。
                これと仮定$F_{A,l+1}(\lambda) = F_{A,l}(\lambda)$より$(A-\lambda I)\bm{v} \in F_{A,l}(\lambda)$。
                よって$(A-\lambda I)^{l+1}\bm{v}=\bm{0}$となり，矛盾する。
            \end{proof}
        %\subsection{相異なる固有値に関する一般固有空間同士の共通部分は$\{\bm{0}\}$}
        %    \begin{shadebox}
        %        $\lambda_i \neq \lambda_j$を$A\in\complexNumbers^{n\times n}$の異なる固有値とし，その代数的重複度を$n_i,n_j$とする。
        %        一般固有空間を$F_i\coloneq\Ker{(A-\lambda_i)^{n_i}},\;F_j\coloneq\Ker{(A-\lambda_j)^{n_j}}$と表すと，$F_i\cap F_j=\{\bm{0}_n\}$
        %    \end{shadebox}
        \subsection{一般固有空間への直和分解}
            \label{一般固有空間への分解}
            \begin{shadebox}
                $A\in\field^{n\times n}$の固有多項式が相異なる固有値$\lambda_1,\dotsc,\lambda_p$に対して$\varphi_A(\lambda) = \prod_{i=1}^{p}{(\lambda-\lambda_i)^{\nu_i}}$となるとき，次が成り立つ。
                \[\field^n = \bigoplus_{i=1}^{p}{F_A(\lambda_i)},\quad \dim{F_A(\lambda_i)} = \nu_i\]
            \end{shadebox}
            \begin{proof}
                \quad\par
                Cayley-Hamiltonの定理と\cite{システム線形代数}補題2.2より最初の式\;$\field^n = \bigoplus_{i=1}^{p}{F_A(\lambda_i)}$\;が従う。
                次に${\nu_i}' \coloneq \dim{F_A(\lambda_i)}$が$\nu_i$に等しいことを示す。
                各$F_A(\lambda_i)$が$A$の不変部分空間であること(\cite{システム線形代数}定理2.4)と定理\ref{不変部分空間と表現行列}より，
                $F_A(\lambda_i)$の基底を$\{\bm{\phi}_{i1},\dotsc,\bm{\phi}_{i{\nu_i}'}\}$とし，$\Phi_i \coloneq [\bm{\phi}_{i1},\dotsc,\bm{\phi}_{i{\nu_i}'}],\quad \Phi \coloneq [\Phi_i,\dotsc,\Phi_p]\in\field^{n\times n}$とすれば，
                次式を満たす正方行列$A_i\in\field^{{\mu_i}'\times{\mu_i}'}\;(i=1,\dotsc,p)$が存在する。
                \[A = \Phi\;\diag{A_1,\dotsc,A_p}\Phi^{-1}\]
                \par
                ここで下準備として$A_i$が唯一の固有値$\lambda_i$を持つことを示しておく。
                まず固有値$\lambda_i$を持つことを示す。$\bm{x}\in E_A(\lambda_i) \subseteq F_A(\lambda_i)$を任意に取り，$F_A(\lambda_i)$の基底で次のように展開する。
                \[\bm{x} = \sum_{j=1}^{{\nu_i}'}{c_j\bm{\phi}_{ij}}\]
                これを上式に右から掛けて
                \begin{align}
                    A\bm{x} &= \Phi\;\diag{A_1,\dotsc,A_p}\Phi^{-1}\sum_{j=1}^{{\nu_i}'}{c_j\bm{\phi}_{ij}} = \Phi\;\diag{A_1,\dotsc,A_p}\sum_{j=1}^{{\nu_i}'}{c_j\bm{e}_{j+\sum_{k=1}^{i-1}{{\nu_k}'}}}\quad(\because \mathrm{定理\ref{逆行列にその構成縦ベクトルを掛ける}}) \nonumber\\
                    &= \Phi\;\diag{A_1,\dotsc,A_p}[\underbrace{0,\dotsc,0}_{{\nu_1}'+\dotsb+{\nu_{i-1}}'},c_1,\dotsc,c_{{\nu_i}'},\underbrace{0,\dotsc,0}_{{\nu_{i+1}}'+\dotsb+{\nu_p}'}]^\top \nonumber\\
                    &= \Phi\left[
                            \begin{array}{c}
                                \bm{0}_{{\nu_1}'}\\
                                \vdots\\
                                \bm{0}_{{\nu_{i-1}}'}\\
                                A_i[c_1,\dotsc,c_{{\nu_i}}']^\top\\
                                \bm{0}_{{\nu_{i+1}}'}\\
                                \vdots\\
                                \bm{0}_{{\nu_p}'}\\
                            \end{array}
                        \right]\quad (\mathrm{※}\bm{0}_\nu\mathrm{は}\nu\mathrm{次元の0ベクトル})
                \end{align}
                ここで$[d_1,\dotsc,d_{\nu_i}']^\top = A_i[c_1,\dotsc,c_{{\nu_i}'}]^\top$とおくと
                \[A\bm{x} = (1) = \sum_{j=1}^{{\nu_i}'}{d_j\bm{\phi}_{ij}}\]
                であるが，そもそも$A\bm{x} = \lambda_i\bm{x} = \sum_{j=1}^{{\nu_i}'}{\lambda_ic_j\bm{\phi}_{ij}}$であったから$d_j = \lambda_jc_j$となり，
                $A_i[c_1,\dotsc,c_{{\nu_i}'}]^\top = \lambda_i[c_1,\dotsc,c_{{\nu_i}'}]^\top$となるから，$A_i$は固有値$\lambda_i$を持つことになる。
                \par
                次に$A_i$の固有値が$\lambda_i$のみであることを背理法で示す。
                $A_i$が固有値$\lambda'\notin\{\lambda_1,\dotsc,\lambda_p\}$を持つと仮定し，$\lambda'$に対する固有ベクトルの一つを$\bm{y} = [y_1,\dotsc,y_{{\nu_i}'}]^\top\in\field^{{\nu_i}'}$とする($A_i\bm{y}=\lambda'\bm{y}$)。
                ベクトル$\bm{x}\in\field^n$を
                \[\Phi^{-1}\bm{x} = [\underbrace{0,\dotsc,0}_{{\nu_1}'+\dotsb+{\nu_{i-1}}'},\bm{y}^\top,\underbrace{0,\dotsc,0}_{{\nu_{i+1}}'+\dotsb+{\nu_p}'}]^\top\]
                すなわち
                \[\bm{x} = \Phi[\underbrace{0,\dotsc,0}_{{\nu_1}'+\dotsb+{\nu_{i-1}}'},\bm{y}^\top,\underbrace{0,\dotsc,0}_{{\nu_{i+1}}'+\dotsb+{\nu_p}'}]^\top = \sum_{j=1}^{{\nu_i}'}{y_j\bm{\phi}_{ij}}\]
                で定めると$\bm{\phi}_{i1},\dotsc,\bm{\phi}_{i{\nu_i}'}$の和で表されているから$\bm{x}\in F_A(\lambda_i)$であるので，まず次式が成り立つ。
                \begin{equation}
                    (A-\lambda_iI_n)^{\nu_i}\bm{x} = \bm{0}_n
                \end{equation}
                一方で
                \begin{align*}
                    (A-\lambda_iI_n)\bm{x} &= A\bm{x} - \lambda_i\bm{x} = \Phi\;\diag{A_1,\dotsc,A_p}\Phi^{-1}\bm{x} - \lambda_i\bm{x}\\
                    &= \Phi\;\diag{A_1,\dotsc,A_p}[\underbrace{0,\dotsc,0}_{{\nu_1}'+\dotsb+{\nu_{i-1}}'},\bm{y}^\top,\underbrace{0,\dotsc,0}_{{\nu_{i+1}}'+\dotsb+{\nu_p}'}]^\top - \lambda_i\bm{x}\\
                    &= \Phi\left[
                        \begin{array}{c}
                            \bm{0}_{{\nu_1}'}\\
                            \vdots\\
                            \bm{0}_{{\nu_{i-1}}'}\\
                            A_i\bm{y}\\
                            \bm{0}_{{\nu_{i+1}}'}\\
                            \vdots\\
                            \bm{0}_{{\nu_p}'}\\
                        \end{array}
                    \right] - \lambda_i\bm{x} = \Phi\left[
                        \begin{array}{c}
                            \bm{0}_{{\nu_1}'}\\
                            \vdots\\
                            \bm{0}_{{\nu_{i-1}}'}\\
                            \lambda'\bm{y}\\
                            \bm{0}_{{\nu_{i+1}}'}\\
                            \vdots\\
                            \bm{0}_{{\nu_p}'}\\
                        \end{array}
                    \right] - \lambda_i\bm{x}\\
                    &= \lambda'\Phi[\underbrace{0,\dotsc,0}_{{\nu_1}'+\dotsb+{\nu_{i-1}}'},\bm{y}^\top,\underbrace{0,\dotsc,0}_{{\nu_{i+1}}'+\dotsb+{\nu_p}'}]^\top - \lambda_i\bm{x} = \lambda'\bm{x}-\lambda_i\bm{x}\\
                    &= (\lambda'-\lambda_i)\bm{x}
                \end{align*}
                であるから
                \[(A-\lambda_iI_n)^{\nu_i}\bm{x} = (\lambda'-\lambda_i)^{\nu_i}\bm{x} \neq \bm{0}_n\]
                となるが，これは式(2)に矛盾する。故に背理法の仮定が誤っており，$A_i$は固有値$\lambda'$を持たない。
                \par
                以上より$A_i$は唯一の固有値$\lambda_i$を持つことがわかる。
                よって
                \[\varphi_A(\lambda) = |\lambda I_n - A| = \prod_{i=1}^{p}{|\lambda I_{{\nu_i}'} - A_i|} = \prod_{i=1}^{p}{(\lambda-\lambda_i)^{{\nu_i}'}}\]
                となり，定理の仮定の$\varphi_A(\lambda)$と比較して${\nu_i}'=\nu_i$が得られる。
            \end{proof}
        この定理から直ちに次が成り立つ。
        \subsection{代数的重複度\texorpdfstring{$\geq$}{≧}幾何学的重複度}
            \begin{shadebox}
                任意の$A\in\field^{n\times n}$の固有値$\lambda$の代数的重複度$\nu_i$は幾何学的重複度$\mu_i$以上である。
            \end{shadebox}
            \begin{proof}
                直前の定理より$\nu_i = \dim{F_A(\lambda)}$であることと，$F_A(\lambda) \supseteq E_A(\lambda)$であることより従う。
            \end{proof}
        \subsection{\texorpdfstring{$AB=O$}{AB=O}でも\texorpdfstring{$\realNumbers^n = \Ker{A} + \Ker{B}$}{R\string^n = Ker(A) + Ker(B)}とは限らない}
            \begin{proof}
                例えば
                $
                    A=B=
                    \begin{bmatrix}
                        0 & 1 \\
                        0 & 0
                    \end{bmatrix}
                $
            \end{proof}
    \section{Jordan標準形}{
        \newcommand{\JBn}{%
            \begin{bmatrix}%
                \lambda & 1          & 0          & \cdots      & 0\\%
                0       & \lambda    & 1          &             & \vdots\\%
                \vdots  & \ddots     & \ddots     & \ddots      & 0\\%
                        &            & 0          & \lambda     & 1\\%
                0       & \cdots     &            & 0           & \lambda%
            \end{bmatrix}%
        }
        \subsection{Jordanブロックの逆行列}
            \begin{shadebox}
                $n$次のJordanブロック
                \[J(\lambda, n) = \JBn\]
                に対して
                \[
                    J(\lambda,n)^{-1} = \left[\indII{i\leq j}\frac{(-1)^{j-i}}{\lambda^{j-i+1}}\right]_{n\times n} =
                    \begin{bmatrix}
                        \frac{1}{\lambda} & \frac{-1}{\lambda^2} & \frac{1}{\lambda^3} & \cdots & \frac{(-1)^{n-1}}{\lambda^n} \\
                        0 & \frac{1}{\lambda} & \frac{-1}{\lambda^2} & \cdots & \frac{(-1)^{n-2}}{\lambda^{n-1}} \\
                        \vdots & 0 & \frac{1}{\lambda} & \ddots & \vdots \\
                        0 & & \ddots & \ddots & \frac{-1}{\lambda^2} \\
                        0 & 0 & \cdots & 0 & \frac{1}{\lambda}
                    \end{bmatrix}
                \]
                \end{shadebox}
                \begin{proof}
                    \quad\newline
                    (M演算(\cite{現代線形代数入門} lesson10)を使う方法):
                    \[ J(\lambda,n)^{-1} = M(\lambda)^{-1} = M(1/\lambda) = \left[\indII{i\leq j}\frac{1}{(i-j)!}\derivLong{\frac{1}{\lambda}}{\lambda}{i-j}\right]_{n\times n} = \left[\indII{i\leq j}\frac{(-1)^{j-i}}{\lambda^{j-i+1}}\right]_{n\times n} \]
                    \newline
                    (M演算を使わない方法):
                    \par
                    帰納法で示す。n次のJordanブロックに対して定理が成り立つと仮定して$(n+1)$次のときも成り立つことを示す。
                    $J(\lambda,n+1)$を次のように分割して
                    \[
                        J(\lambda,n+1) =
                        \begin{bmatrix}
                            J(\lambda,n) & \bm{b} \\
                            \bm{0}_n^\top & \lambda
                        \end{bmatrix}
                        ,\quad
                        \bm{b} = [0\cdots0\;1]^\top \in \complexNumbers^n
                    \]
                    と表すと，これの逆行列は次のように表される(\cite{入門線形代数} 問題2.4.8)。
                    \[
                        J(\lambda,n+1)^{-1} =
                        \begin{bmatrix}
                            J(\lambda,n)^{-1} & \bm{v} \\
                            \bm{0}_n^\top & 1/\lambda
                        \end{bmatrix}
                        ,\quad
                        \bm{v} = \frac{-1}{\lambda}J(\lambda,n)^{-1}\bm{b} = \left[\frac{-1}{\lambda}J(\lambda,n)^{-1}[i][n]\right]_{n\times 1} = \left[\frac{(-1)^{(n+1)-i}}{\lambda^{(n+1)-i+1}}\right]_{n\times 1}
                    \]
                    よって$n+1$次のときも定理が成り立つ。
                \end{proof}
        \subsection{Jordanブロックの指数行列}
            \begin{shadebox}
                $n$次のJordanブロック
                \[J(\lambda, n) = \JBn\]
                に対して
                \[
                    \NapierE^{J(\lambda, n)t} =
                    \begin{bmatrix}
                        \NapierE^{\lambda t}    & t\NapierE^{\lambda t}    & \cdots    & \frac{t^{n-1}}{(n-1)!}\NapierE^{\lambda t}\\
                        0                & \NapierE^{\lambda t}        & \ddots    & \vdots\\
                        \vdots            & \ddots            & \ddots    & t\NapierE^{\lambda t}\\
                        0                & \cdots            & 0            & \NapierE^{\lambda t}
                    \end{bmatrix}
                \]
            \end{shadebox}
            \begin{proof}
                \quad\par
                $J(\lambda, n)$を次のように対角成分とその一つ上の斜めラインに分ける。
                \[
                    J(\lambda, n)=\Gamma+U,\quad
                    \Gamma\coloneq\lambda I_{n\times n},\quad
                    U\coloneq
                    \begin{bmatrix}
                        0        & 1    & 0            & \cdots    & 0\\
                                & 0    & 1            & \ddots    & \vdots\\
                        \vdots    &    & \ddots    & \ddots    & 0\\
                                &     &             & 0            & 1\\
                        0        &     & \cdots    &             & 0
                    \end{bmatrix}
                \]
                \par
                $\Gamma$と$U$は可換であり($\Gamma$が単位行列の定数倍であることから明らか)，$\Gamma$に$U$を左或いは右から掛けると$\Gamma$の各行が上に一段シフトする(或いは各列が右に1段シフトするとも言える)。
                溢れた行或いは列は消え去り，補充されない。溢れた分だけ反対側から空行或いは空列が注入される。
                よって$\left[J(\lambda, n)t\right]^p,\quad p\in\mathbb{N}+\{0\}$の対角成分より下側の成分は$0$であることがわかる。
                \par
                \begin{equation}
                    \left[J(\lambda, n)t\right]^p = t^p(\Gamma+U)^p = t^p\sum_{k=0}^{p}{\combi{p}{k}U^k\Gamma^{p-k}} \tag{1}
                \end{equation}
                $(\Gamma+U)^p$の$(i,j)$成分を$a_{i,j}(p)$とする。先述の通り$a_{i,j}(p)=0,\;i>j$である。
                \par
                以下$i\leq j$とする。式(1)の$U^k\Gamma^{p-k}$において$(i,j)$要素が非零となるのは$k=j-i$のときに限る(当然，同時に$p\geq j-i$でなくてはならない)。
                各$p$に対する$a_{i,j}(p)$を表にすると次のようになる。
                \begin{table}[H]
                    \centering
                    \caption{$p$と$a_{i,j}(p)$}
                    \begin{tabular}{cc}
                        $p$            & $a_{i,j}(p)$ \\\hline\hline
                        $0$            & $0$\\
                        $1$            & $0$\\
                        $\vdots$    & $\vdots$\\
                        $j-i$        & $\combi{p}{j-i}$\\
                        $j-i+1$        & $\lambda\combi{p}{j-i}$\\
                        $j-i+2$        & $\lambda^2\combi{p}{j-i}$\\
                        $\vdots$    & $\vdots$\\
                        $p$            & $\lambda^{p-(j-i)}\combi{p}{j-i}$\\\hline
                    \end{tabular}
                \end{table}
                よって$\NapierE^{J(\lambda, n)t}$の$(i,j)$成分は
                \begin{align*}
                    \sum_{p=0}^{\infty}{\frac{1}{p!}t^p a_{i,j}(p)} &= \sum_{p=j-i}^{\infty}{\frac{1}{p!}t^p\lambda^{p-(j-i)}\combi{p}{j-i}}\\
                    &= \sum_{l=0}^{\infty}{\frac{1}{\left(l+(j-i)\right)!}t^{l+(j-i)}\lambda^l\frac{\left(l+(j-i)\right)!}{(j-i)!l!}}\quad(l=p-(j-i)\mathrm{とおいた})\\
                    &= \frac{t^{j-i}}{(j-i)!}\sum_{l=0}^{\infty}{\frac{(t\lambda)^l}{l!}} = \frac{t^{j-i}}{(j-i)!}\NapierE^{\lambda t}
                \end{align*}
            \end{proof}
    }
    \section{Perronの定理}
        \begin{shadebox}
            $n$次正行列$A$について次が成り立つ
            \begin{enumerate}
                \item $\rho(A)>0$
                \item $\rho(A)$は$A$の単純固有値の一つと等しく，$\rho(A)$は他のどの固有値の絶対値よりも真に大きい
                \item $\rho(A)$と等しい固有値に対応する正の固有ベクトルが存在する
            \end{enumerate}
        \end{shadebox}
        Perron-Frobeniusの定理の主張との違いは 2. の「$\rho(A)$は他のどの固有値の絶対値よりも真に大きい」である。
        \begin{proof}
            \quad\par
            1.と3.および2.の「$\rho(A)$は$A$の単純固有値の1つと等しく」は\cite{システム線形代数}のPerron-Frobeniusの定理の証明と同じ方法で証明できる。
            2.の「$\rho(A)$は他のどの固有値の絶対値よりも真に大きい」を，\cite{システム線形代数} 170頁 定理6.11 のPerron-Frobeniusの定理の証明を利用して示す。
            その証明の中に
            \begin{equation}
                \lambda_ix_i = \sum_{j=1}^{n}{a_{ij}x_j}
            \end{equation}
            \[|\lambda||x_i| = \left|\sum_{j=1}^{n}{a_{ij}x_j}\right| \leq \sum_{j=1}^{n}{a_{ij}|x_j|}\]
            という式があるが，$|\lambda|\leq\hat{\lambda}$は既に示されているので，上式に於いて等号が成り立てば$\lambda=\hat{\lambda}$であることを示せばよい。
            \par
            $A$は正行列だから$a_{ij}>0$である。$\bm{x}$は固有ベクトルであるから$x_j$の少なくともどれか一つは非零である。
            よって等号が成り立つとき$x_j$は全て同じ向きである。
            言い換えれば，ある非零な複素数$z$と$b_1,\dotsc,b_n\geq0$が存在して
            \[
                \bm{x} = z
                \begin{bmatrix}
                    b_1\\
                    b_2\\
                    \vdots\\
                    b_n
                \end{bmatrix}
            \]
            である。これを式(1)に代入して
            \begin{align*}
                \lambda zb_i &= \sum_{j=1}^{n}{a_{ij}zb_j}\\
                \therefore\; \lambda b_i &= \sum_{j=1}^{n}{a_{ij}b_j}
            \end{align*}
            $b_1,\dotsc,b_n$の少なくとも1つは正である。それを$b_k$とすると
            \[\lambda = \frac{1}{b_k}\sum_{j=1}^{n}{a_{ij}b_j} > 0\]
            これと$|\lambda|=\hat{\lambda}$より$\lambda=\hat{\lambda}$
        \end{proof}